% Optional: Customize the appearance of the code (Set for a decent MATLAB code style)
\lstset{
    language=Matlab,           % Set language to MATLAB
    basicstyle=\ttfamily\small, % Basic style for the code font
    keywordstyle=\color{blue},  % Color for keywords
    commentstyle=\color{green!70!black}, % Color for comments
    stringstyle=\color{violet!50!white},    % Color for strings
    numbers=left,               % Line numbers on the left
    numberstyle=\tiny\color{gray}, % Style for line numbers
    stepnumber=1,               % Line number step
    frame=single,               % Frame around the code
    breaklines=true,            % Automatically break long lines
    captionpos=b,               % Caption position at the bottom
    showspaces=false,           % Do not show spaces
    showstringspaces=false,     % Do not show string spaces
    literate={Ω}{{$\Omega$}}1  % Replace Ω with LaTeX \textOmega
             {Δ}{{$\Delta$}}1      % Replace Δ with LaTeX \Delta
             {°}{{\degree}}1     % Replace ° with LaTeX \degree
}

% Define G-code as a custom language
\lstdefinelanguage{GCode}{
  morekeywords={G0,G1,G2,G3,G4,G21,G28,G90,G91,M0,M1,M2,M3,M4,M5,M6,M30,M82,M83,F,S,T,X,Y,Z,E}, 
  sensitive=false,
  morecomment=[l];,
  morestring=[b]"
}

% Define the colors for clarity
\definecolor{codesysblue}{rgb}{0.0, 0.0, 1}
\definecolor{codesysgreen}{rgb}{0.0, 0.5, 0.0}
\definecolor{codesysmagenta}{rgb}{0.8, 0.0, 0.8}

\lstdefinelanguage{CoDeSysV23}{
    keywords={VAR_GLOBAL,END_VAR,VAR,END_VAR,AT,BOOL,INT,REAL,STRING,TIMER,TON,TOF,CTU,IF,THEN,ELSE,ELSIF,END_IF,FOR,TO,DO,END_FOR,WHILE,END_WHILE,REPEAT,UNTIL,END_REPEAT,CASE,OF,END_CASE,EXIT,RETURN,NOT,AND,OR,XOR,:=,FUNCTION_BLOCK,VAR_INPUT,VAR_OUTPUT},
    % TRUE/FALSE are now listed as morekeywords, and their style is customized.
    morekeywords=[2]{TRUE,FALSE}, 
    sensitive=false,
    morecomment=[l]{(*},
    morecomment=[l]{//},
    commentstyle=\color{codesysgreen},
    keywordstyle=\color{codesysblue},
    % Define style for second set of keywords (TRUE/FALSE)
    keywordstyle={[2]\color{codesysmagenta}}, 
    stringstyle=\color{red},
    identifierstyle=\color{black},
    % Now we use the style defined above, and focus literate only on addresses.
    literate= 
    % NEW: Highlight T# constant prefix in pink ---
             {T\#}{{{\color{codesysmagenta} T\#}}}2 % T# has 2 characters
    % --- Pink Elements (Addressess + Space) ---
             {\%IX2.0\ }{{{\color{codesysmagenta} \%IX2.0}}}6
             {\%IX2.1\ }{{{\color{codesysmagenta} \%IX2.1}}}6
             {\%IX2.2\ }{{{\color{codesysmagenta} \%IX2.2}}}6
             {\%IX2.3\ }{{{\color{codesysmagenta} \%IX2.3}}}6
             {\%IX2.4\ }{{{\color{codesysmagenta} \%IX2.4}}}6
             {\%IX2.5\ }{{{\color{codesysmagenta} \%IX2.5}}}6
             {\%IX2.6\ }{{{\color{codesysmagenta} \%IX2.6}}}6
             {\%IX2.7\ }{{{\color{codesysmagenta} \%IX2.7}}}6
             {\%IX4.0\ }{{{\color{codesysmagenta} \%IX4.0}}}6
             {\%IX4.1\ }{{{\color{codesysmagenta} \%IX4.1}}}6
             {\%IX4.2\ }{{{\color{codesysmagenta} \%IX4.2}}}6
             {\%IX4.3\ }{{{\color{codesysmagenta} \%IX4.3}}}6
             {\%IX4.4\ }{{{\color{codesysmagenta} \%IX4.4}}}6
             {\%IX4.5\ }{{{\color{codesysmagenta} \%IX4.5}}}6
             {\%IX4.6\ }{{{\color{codesysmagenta} \%IX4.6}}}6
             {\%IX4.7\ }{{{\color{codesysmagenta} \%IX4.7}}}6
             {\%QX0.0\ }{{{\color{codesysmagenta} \%QX0.0}}}6
             {\%QX0.1\ }{{{\color{codesysmagenta} \%QX0.1}}}6
             {\%QX0.2\ }{{{\color{codesysmagenta} \%QX0.2}}}6
             {\%QX0.3\ }{{{\color{codesysmagenta} \%QX0.3}}}6
             {\%QX0.4\ }{{{\color{codesysmagenta} \%QX0.4}}}6
             {\%QX0.5\ }{{{\color{codesysmagenta} \%QX0.5}}}6
             {\%QX0.6\ }{{{\color{codesysmagenta} \%QX0.6}}}6
             {\%QX2.0\ }{{{\color{codesysmagenta} \%QX2.0}}}6
             {\%QX2.1\ }{{{\color{codesysmagenta} \%QX2.1}}}6
             {\%QX2.2\ }{{{\color{codesysmagenta} \%QX2.2}}}6
             {\%QX2.3\ }{{{\color{codesysmagenta} \%QX2.3}}}6
             {\%QX2.4\ }{{{\color{codesysmagenta} \%QX2.4}}}6
             {\%QX2.5\ }{{{\color{codesysmagenta} \%QX2.5}}}6
             {\%QX2.6\ }{{{\color{codesysmagenta} \%QX2.6}}}6
             {\%QX2.7\ }{{{\color{codesysmagenta} \%QX2.7}}}6
             {T\#250ms}{{{\color{codesysmagenta} T\#250ms}}}7
             {T\#300ms}{{{\color{codesysmagenta} T\#300ms}}}7
             {T\#440ms}{{{\color{codesysmagenta} T\#440ms}}}7
             {T\#200ms}{{{\color{codesysmagenta} T\#200ms}}}7
             {T\#10ms}{{{\color{codesysmagenta} T\#10ms}}}7
             {T\#100ms}{{{\color{codesysmagenta} T\#100ms}}}7,
}
\makeatother % End of internal LaTeX command usage

% Define the style based on the new language
\lstdefinestyle{codesys}{
    language=CoDeSysV23, % Use the newly defined language
    basicstyle=\footnotesize\ttfamily,
    numbers=left,
    numberstyle=\tiny,
    stepnumber=1,
    numbersep=8pt,
    showspaces=false,
    showtabs=false,
    breaklines=true,
    frame=single,
    frameshape={l}{}{}{},
    framerule=0.4pt,
    tabsize=4,
    % No need to redefine styles here, as they are set in the language definition
}